\chapter{RPA 文本汇编与工具链}
\label{ch:assembly}

RPA（Reformulation Proof Assembly）是 RPX 的文本表示。文本层承担符号解析、精确数值解析、表达式展开和二进制编码；证书是否成立仍由 RPX verifier 判定。本章定义用于人工编写、生成器输出和反汇编的基础语法，不赋予 assembler 任何独立的证明权力。

本章中的命令行和工具名称是 v0.1 工具接口草案。它们描述计划交付的行为，不表示当前仓库已经实现这些程序。MCM 的二进制编码和完整 MPS 导入覆盖范围尚待单独冻结。

\section{文件结构}

\subsection{编码与词法约定}

RPA 文件采用 UTF-8 编码，不含 BOM。语言关键字、标识符和操作数字符使用 ASCII；注释可以使用 UTF-8。关键字区分大小写，规范输出采用本章给出的小写拼写。空格和制表符是 token 分隔符；行结束可以是 LF 或 CRLF。assembler 在词法处理时将 CRLF 视为 LF。

\code{;} 引入行注释，注释持续到该行结束。基础语法不支持块注释、宏、条件汇编、文件包含或用户定义伪指令。分号不是语句分隔符。除花括号中的列表外，一条语句必须在同一物理行结束；列表允许跨行。文件末尾可以省略最终换行。

标识符由字母或下划线起始，随后可以包含字母、数字和下划线。Fact 和 expression 标签另外允许全数字名称。这里的数字名称仍然是\emph{符号名称}，不是用户指定的二进制 ID。例如作者写的 \code{\%99} 可以被分配为第一个 derived FactId；规范反汇编则写成 \code{\%0}。同一 execution domain 内，同类标签不得重复定义。两个 domain 的标签命名空间独立。

源码中的变量、行和静态事实索引必须是非负十进制整数。前导零允许被 assembler 接受，规范输出必须删除前导零。引用索引必须满足对应模型或 FactRef 编码的取值范围；表达式不得引用输出变量。

\subsection{顶层声明}

一个基础 v0.1 文件包含版本、两个模型摘要、contract 和两个 execution domain。\code{.forward} 和 \code{.reverse} 各出现一次，并按此顺序排列。每个 domain 恰好包含一个 \code{.map} 区域和一个 \code{.code} 区域；所有 \code{.expr} 声明位于 \code{.map} 之前。

\begin{lstlisting}[caption={RPA 文件结构；尖括号表示说明性占位符}]
.rpa 0.1
.source sha256 "<64 hexadecimal digits>"
.target sha256 "<64 hexadecimal digits>"
.contract mip=opt-equivalent
.contract forward.obj=eq
.contract reverse.obj=eq

.forward
.expr $name = <affine expression>
.map
    <static mapping declarations>
.endmap
.code
    <proof instructions>
    ret
.endcode

.reverse
.expr $other = <affine expression>
.map
    <static mapping declarations>
.endmap
.code
    <proof instructions>
    ret
.endcode
\end{lstlisting}

\code{.source} 和 \code{.target} 的字符串必须各包含恰好 64 个十六进制数字。输入可以使用大小写字母，规范输出使用小写。摘要是相应 canonical model 的 SHA-256，不是原始 MPS 文本的摘要。hash 字符串不支持转义或空白。

\begin{longtable}{p{0.32\textwidth}p{0.60\textwidth}}
\toprule
声明 & 基础 v0.1 取值与约束 \\
\midrule
\endhead
\code{.rpa 0.1} & 必须为第一条非注释语句；一个文件只出现一次。\\
\code{.contract mip=opt-equivalent} & 要求双向可行映射与双向目标关系；每个方向的目标关系必须为 \code{eq} 或 \code{le}。\\
\code{.contract mip=feasibility-equivalent} & 要求双向可行映射；允许 \code{none}、\code{eq} 或 \code{le}。指定 EQ/LE 时还必须完成对应目标义务。\\
\code{.contract forward.obj=eq} & Forward 的归一化目标差要求等于零。\\
\code{.contract forward.obj=le} & Forward 的归一化目标差要求不大于零。\\
\code{.contract forward.obj=none} & 仅用于可行性等价 contract；不生成 Forward objective obligation。\\
\code{.contract reverse.obj=...} & 与 Forward 相同，但 input/output 模型方向交换。\\
\bottomrule
\end{longtable}

每条 contract 声明必须恰好出现一次，不能通过重复声明覆盖先前取值。Forward 和 Reverse 共同证明文件声明的 contract。调用者要求的最小 contract 由 verifier 的调用参数决定；文件不能自行降低调用者要求。最大化目标和最小化目标的归一化定义见机器模型与 contract 章节。

\section{精确数值字面量}

\subsection{允许的拼写}

基础语法接受整数、分数、有限十进制数和十进制科学计数法。分数字面量的分母必须为严格正整数；符号只放在分子之前。

\begin{lstlisting}
17
-3
4/7
-22/13
0.125
.125
1.
1e-6
-2.5E+3
\end{lstlisting}

这些写法均通过整数运算解析，不经过 IEEE 浮点数。十进制和科学计数法定义如下。将十进制尾数的全部数字拼接为整数 $d$，设小数点后有 $k$ 位，指数为整数 $t$，符号为 $s\in\{-1,1\}$，则：
\[
  \operatorname{value}=s\,d\,10^{t-k}.
\]
例如 \code{0.125} 解析为 $125/1000=1/8$，\code{1e-6} 解析为 $1/1000000$，\code{-2.5E+3} 解析为 $-2500$。分数通过约分得到 canonical rational。\code{-0}、\code{0/7} 和 \code{0.0} 均归一化为整数零。

\must{assembler 必须以精确整数运算实现数值字面量转换，不能先解析为浮点数再尝试恢复有理数。}

基础语法不接受 \code{NaN}、\code{Inf}、十六进制数、数字分组下划线、循环小数或隐含分母符号。\code{1/0} 是错误；\code{1/-2} 必须写为 \code{-1/2}。资源限额可以限制字面量长度和指数绝对值，但不得改变其数学值。

\subsection{规范输出}

规范反汇编只输出整数或约分后的 \code{p/q}：分母严格正，整数不带 \code{/1}，正数不带前导加号，零写作 \code{0}。因此源码中的十进制拼写、科学计数法、前导零和原始分数形式不会被保留。

\section{操作数与符号}

\begin{longtable}{p{0.29\textwidth}p{0.28\textwidth}p{0.35\textwidth}}
\toprule
操作数形式 & 类别 & 意义 \\
\midrule
\endhead
\code{\%name} & Derived FactRef & 当前 domain 已定义且可见的事实。\\
\code{@in.r[n].hi} & 静态 FactRef & input row 的有限 upper side，类型为 LE。\\
\code{@in.r[n].lo} & 静态 FactRef & input row 的有限 lower side，类型为 LE。\\
\code{@in.r[n].eq} & 静态 FactRef & input equality row，类型为 EQ。\\
\code{@in.v[n].hi} & 静态 FactRef & input variable 的有限 upper bound，类型为 LE。\\
\code{@in.v[n].lo} & 静态 FactRef & input variable 的有限 lower bound，类型为 LE。\\
\code{@in.v[n].int} & 静态 FactRef & input variable 的 integrality，类型为 INT。\\
\code{@case} & Branch FactRef & 当前最内层分支假设；等价于 \code{@case[0]}。\\
\code{@case[n]} & Branch FactRef & 距当前分支 $n$ 层的祖先分支假设。\\
\code{\$name} & ExprId 符号 & 当前 domain 的 canonical affine expression。\\
\code{@in.v[n]} & Input VarId & 表达式或 mapping 中的 input variable。\\
\code{@out.v[n]} & Output VarId & mapping 或 variable obligation 的对象。\\
\code{@out.r[n]} & Output RowId & row obligation 的对象。\\
\bottomrule
\end{longtable}

静态 FactRef 不产生指令，也不占用 derived FactId。只有模型确实声明的事实可以被引用：没有有限 upper bound 时，\code{@in.v[n].hi} 不存在；连续变量没有 \code{.int} 事实；非 equality row 没有 \code{.eq} 事实。equality row 的 side 可用性遵循 canonical model 对有限 sides 的定义。

\code{@case[n]} 中的 $n$ 是 lexical scope depth offset，不是分支编号。当不存在该层假设时引用无效。\code{@case} 的具体类型和表达式由 split 指令规定，不能由源码重声明或赋值。

\section{表达式声明与展开}

\subsection{仿射语法}

表达式语法支持精确常数、input variable、此前声明的 expression 符号、括号、加减法及标量乘法：

\begin{lstlisting}
.expr $x = @in.v[0]
.expr $sum = @in.v[0] + @in.v[2] - 1
.expr $scaled = 3/5 * $sum + 2
.expr $neg = -$scaled
.expr $mixed = 2 * (@in.v[0] - @in.v[2]) + 7
\end{lstlisting}

expression 符号通过展开得到 affine expression；符号不是运行时节点。引用必须指向同一 domain 内此前声明的 expression，因此展开关系无环。每个声明最终归一化为一个常数和严格按 VarId 排序的非零系数列表。

不允许 variable 与 variable 相乘、不允许两个一般表达式相乘，也不允许除以一般表达式。乘法的左侧必须是一个数值字面量。\code{3 * @in.v[0]} 合法；\code{@in.v[0] * 3} 不属于基础语法。减法、负号和括号可用于任意 affine expression。\code{2 @in.v[0]} 这样的隐含乘法不合法。

\subsection{指令中的 expression 操作数}

基础指令的 expression 操作数必须为 \code{\$name}。例如 \code{map.aff @out.v[0], -\$e} 不合法，必须先声明 \code{.expr \$neg = -\$e}，再使用 \code{map.aff @out.v[0], \$neg}。这一限制确保 instruction 编码直接引用 expression pool，并让反汇编语法唯一。

assembler 可以合并数值相同的 expression 或 scalar pool 条目；它必须保留程序的数学语义和 scope 行为。基础文件格式是否允许重复 pool 条目由 wire format 章节规定，文本标签不承担 pool 唯一性保证。

\section{静态 mapping 区域}

\code{.map} 描述不依赖分支的 output variable mapping。它的变量命名方向随 domain 改变：Forward 中 input 是 Source、output 是 Target；Reverse 中相反。

\begin{lstlisting}
.map
    copy.range @out.v[0], @in.v[0], 1000
    copy @out.v[1000], @in.v[7]
    const @out.v[1001], 0
    aff @out.v[1002], $replacement
    defer @out.v[1003]
    defer.range @out.v[1004], 4
.endmap
\end{lstlisting}

\begin{longtable}{p{0.18\textwidth}p{0.39\textwidth}p{0.35\textwidth}}
\toprule
语句 & 文本操作数 & section 编码与含义 \\
\midrule
\endhead
\code{copy.range} & output 起点、input 起点、count & COPY\_RANGE；连续 $count$ 对变量作 identity-copy。\\
\code{copy} & output variable、input variable & COPY\_PAIRS 的单对声明；assembler 可以合并多条声明。\\
\code{const} & output variable、number & CONST；output 等于该精确常数。\\
\code{aff} & output variable、expression & AFF；output 等于该 input affine expression。\\
\code{defer} & output variable & DEFER；该变量必须在 proof tape 的可达路径上定义。\\
\code{defer.range} & output 起点、count & 连续 DEFER 声明的文本缩写。\\
\bottomrule
\end{longtable}

\must{静态 map 必须恰好覆盖全部 output variables；每个变量恰好由 COPY、CONST、AFF 或 DEFER 覆盖一次。}

range 的 count 必须为严格正整数，索引范围的加法必须检查溢出。两个范围重叠、copy 越界或重复声明均是错误。一个静态 DEFER 标记不是已经完成的 mapping，也不能消除任何 domain obligation。

\code{map.copy}、\code{map.const} 和 \code{map.aff} 是 proof tape 指令，语法与相应静态声明相近。它们只用于当前路径尚未赋值的 DEFER output variable。对已有 mapping 重复赋值不合法。不同 sibling branches 可以各自给同一个 DEFER variable 赋值，因为两个分支分别从父状态执行。

\section{基础指令的文本形式}

下列语法穷尽基础 v0.1 的 28 个 core opcodes。所有指令的数学检查、异常条件和二进制 payload 见 ISA 章节；本节只规定文本 spelling 和操作数顺序。

\subsection{General、linear 与 integer 指令}

\begin{lstlisting}
nop
ret
drop %fact
drop %first .. %last

%out = lin.le $expr {
    +1 @in.r[0].hi
    -3/5 @in.r[1].eq
}
%out = lin.eq $expr {
    +1 @in.r[0].eq
    -2 %earlier
}
%out = eq.pair $expr, %positive, %negative
%out = box.le $expr
%out = weaken.le $expr, %input, 3/2
close.false %contradiction

%out = int.aff $expr, 7 {
    +3 @in.v[0].int
    -2 %integer_fact
    +1/2 %equality_fact
}
%out = round.hi $expr, %upper, %integer
%out = round.lo $expr, %lower, %integer
\end{lstlisting}

线性列表中的每一行恰好包含一个 multiplier 和一个 FactRef，中间至少一个空白，不加逗号。允许空列表；空 \code{lin.eq} 或 \code{lin.le} 只能证明零表达式。\code{int.aff} 的 offset 必须是精确整数。整数事实的 multiplier 必须为整数；EQ 项的 multiplier 可以是任意有理数，这是整数性通过等式传递的草案修订 P-04。

\code{drop} 的区间端点经符号解析后必须形成递增的连续 FactId 区间，其中每个 ID 均为当前 scope 已定义的 live fact。区间包含未来、dead 或其他 scope 的事实时不合法；不能因文本标签恰好具有连续数字而假定其机器 ID 连续。assembler 可以将 liveness 结果拆成多条合法 DROP 区间。

\subsection{Mapping 与 obligation 指令}

\begin{lstlisting}
map.copy @out.v[0], @in.v[1]
map.const @out.v[1], 0
map.aff @out.v[2], $expr

check.row.hi @out.r[0], %fact
check.row.lo @out.r[0], @in.r[1].lo
check.row.eq @out.r[0], %fact
check.var.hi @out.v[0], %fact
check.var.lo @out.v[0], @in.v[1].lo
check.var.int @out.v[0], %fact
check.obj.eq %fact
check.obj.le %fact
\end{lstlisting}

CHECK 的 proof operand 可以是任意合法 FactRef，包括静态模型事实和 branch assumption，无需先复制到 derived register。其类型和表达式必须与 kernel 生成的 obligation 精确一致。

\subsection{\code{exact} 简写}

\proposal{P-05}{为 CHECK 增加 \code{exact} 文本操作数；编译为 CHECK 的 EXACT flag，并由 kernel 直接检查下述有限规则。它不是 assembler 信任断言。}

\begin{lstlisting}
check.row.eq @out.r[0], exact
check.var.int @out.v[0], exact
check.obj.eq exact
\end{lstlisting}

exact 的判断对象始终是 kernel 自己由 output model 和当前 mapping 构造的表达式 $e$：

\begin{longtable}{p{0.18\textwidth}p{0.74\textwidth}}
\toprule
所需 FactKind & EXACT 可接受条件 \\
\midrule
\endhead
EQ & canonical $e$ 为恒等零。\\
LE & canonical $e$ 无变量项，且其常数不大于零。\\
INT & $e$ 的常数和所有系数均为整数；所有非零项涉及的 input variables 在 input model 中声明为整数或 binary。\\
\bottomrule
\end{longtable}

EXACT 不搜索 input rows，不自动匹配 input bounds，也不推导未知变量的整数性。相同 input row 直接用于 output row 时应显式写出对应静态 FactRef。一个必须依赖已验证 INT fact 的整数性检查，应使用该 fact 或 \code{int.aff}，不能要求 EXACT 发现依赖。

模型相关的 \code{exact} 预检查和 expression 归一化要求 assembler 获得两个 canonical models。即使 assembler 已预检查成功，verifier 仍必须独立执行同一判定；源文本不能成为可信输入。

\subsection{Structured cases}

\begin{lstlisting}
(%a, %b) = split.bin @in.v[0] {
case 0:
    yield %left_a, %left_b
case 1:
    yield %right_a, %right_b
}

(%a) = split.int.thr %integer, 13 {
case le:
    yield %left
case ge:
    yield %right
}

(%a) = split.int.range %integer, %lower, %upper, 3, 5 {
case 3:
    yield %value3
case 4:
    yield %value4
case 5:
    yield %value5
}

() = split.sign $expr {
case le:
    yield
case ge:
    yield
}
\end{lstlisting}

split 总是带输出 tuple；零个输出使用 \code{()}。\code{split.bin} 的 cases 必须按 \code{0}、\code{1} 顺序出现；\code{split.int.thr} 和 \code{split.sign} 的 cases 必须按 \code{le}、\code{ge} 顺序出现；\code{split.int.range} 的 cases 必须按闭区间中的整数递增排列。case labels 不是可跳转地址。

所有 sibling branches 必须从同一父状态开始执行。若一个 DEFER variable 在部分可达 branches 中已赋值、另一些可达 branches 中仍未赋值，merge 立即返回 \code{E_MAP}。只有所有可达 branches 均未赋值时才保留 UNASSIGNED；均已赋值时按 mapping merge 规则保留共同 expression 或封存为 SEALED\_PIECEWISE。

每个 case 必须以 \code{yield} 或 \code{close.false} 结束，之后不能有指令。\code{yield} 的 FactRef 数量必须等于 split 输出 tuple 长度；\code{yield} 不带操作数表示零个事实。所有可达 branches 在每个 yield 位置必须产生完全相同的 canonical proposition。已证明不可达的 branches 不需要 yield。split 可以嵌套，但不允许 JMP 或其他非结构化跳转。

\code{split.int.thr} 的 threshold 和 \code{split.int.range} 的两个端点必须为精确整数。对于 INT expression $g$，range 所用 lower fact 必须精确表示 $-g+L\le0$，upper fact 必须精确表示 $g-U\le0$。闭区间端点要求 $L\le U$。

对于 \code{split.sign}，两个 proof branches 的假设分别为 $e\le0$ 和 $-e\le0$，零点属于两支。mapping 的数值应用若遇到零，按草案 P-06 选择先出现的 \code{le} case；验证仍检查两个 proof branches。

\section{FactId、scope 与规范反汇编}

\proposal{P-03}{derived FactId 按语法顺序隐式分配：普通 producer 各分配一个 ID；split 先递归分配各 branch 的 producer IDs，再分配 merge outputs。ID 分配不依赖数学可达性。}

assembler 先解析全部符号和结构，依此分配 FactId；kernel 仍禁止未来事实引用。split 输出在源码头部命名，但只有 split 完成后可见，不能在其 branches 内引用。branch-local facts 不能通过符号泄漏到其他 sibling branch 或父 scope。源代码在文本上较早定义某个 fact，不表示该 fact 在另一个 scope 中可见。

规范反汇编将 expression 符号重命名为 \code{\$e0}、\code{\$e1} 等，将 derived facts 重命名为机器 FactId 的十进制形式。它输出固定 case 顺序、约分数值、展开后的 canonical expressions 和确定的空白规则。

\must{assemble--disassemble--assemble 必须保持 contract、数学 propositions、控制结构和 mapping 语义；基础 v0.1 不承诺逐字节还原原始文件。}

原因包括 scalar/expression pool 的合法重复条目、section 排列、debug 数据以及文本标签可能不同。若需要逐字节保存，工具必须提供独立的原始二进制导出模式，不能把这种需求隐藏在 canonical RPA 语义中。

\section{核心 EBNF}

以下 EBNF 使用 \code{[...]} 表示可选部分、\code{\{...\}} 表示重复、\code{|} 表示选择。\code{NL} 表示换行，\code{WS} 表示至少一个水平空白。注释在语法分析前处理；空白行允许出现在任意完整语句之间。表达式加减运算左结合；scalar 乘法优先于加减法，单目符号优先于加减法。

\begin{lstlisting}[caption={词法与顶层语法}]
digit       = "0" | "1" | ... | "9" ;
letter      = "A" | ... | "Z" | "a" | ... | "z" ;
hex-digit   = digit | "A" | ... | "F" | "a" | ... | "f" ;
WS          = (" " | TAB), { " " | TAB } ;
NL          = LF ;
uint        = digit, { digit } ;
sign        = "+" | "-" ;
integer     = [ sign ], uint ;
fraction    = [ sign ], uint, "/", uint ;
decimal     = uint, ".", [ uint ] | ".", uint ;
mantissa    = uint | decimal ;
scientific  = [ sign ], mantissa, ("e" | "E"), integer ;
number      = fraction | scientific | [ sign ], decimal | integer ;
identifier  = (letter | "_"), { letter | digit | "_" } ;
name        = identifier | uint ;
fact-name   = "%", name ;
expr-name   = "$", name ;
input-var   = "@in.v[", uint, "]" ;
output-var  = "@out.v[", uint, "]" ;
output-row  = "@out.r[", uint, "]" ;
fact-ref    = fact-name
            | "@in.r[", uint, "].", ("hi" | "lo" | "eq")
            | "@in.v[", uint, "].", ("hi" | "lo" | "int")
            | "@case", [ "[", uint, "]" ] ;
hash        = '"', 64 * hex-digit, '"' ;
objective   = "eq" | "le" | "none" ;
contract    = ".contract mip=", ("opt-equivalent"
              | "feasibility-equivalent"), NL,
              ".contract forward.obj=", objective, NL,
              ".contract reverse.obj=", objective, NL ;
file        = ".rpa 0.1", NL,
              ".source sha256", WS, hash, NL,
              ".target sha256", WS, hash, NL,
              contract, forward-domain, reverse-domain ;
forward-domain = ".forward", NL, domain-body ;
reverse-domain = ".reverse", NL, domain-body ;
domain-body = { expression-declaration },
              ".map", NL, { map-declaration }, ".endmap", NL,
              ".code", NL, code-body, ".endcode", [ NL ] ;
\end{lstlisting}

\begin{lstlisting}[caption={仿射表达式与静态 mapping 语法}]
expression-declaration = ".expr", WS, expr-name, "=", affine, NL ;
affine      = scaled, { ("+" | "-"), scaled } ;
scaled      = number, "*", signed-primary | signed-primary ;
signed-primary = [ sign ], primary ;
primary     = number | input-var | expr-name | "(", affine, ")" ;

map-declaration = "copy.range", WS, output-var, ",", input-var,
                    ",", uint, NL
                | "copy", WS, output-var, ",", input-var, NL
                | "const", WS, output-var, ",", number, NL
                | "aff", WS, output-var, ",", expr-name, NL
                | "defer", WS, output-var, NL
                | "defer.range", WS, output-var, ",", uint, NL ;
\end{lstlisting}

\begin{lstlisting}[caption={普通指令语法}]
linear-list = "{", NL, { number, WS, fact-ref, NL }, "}" ;
proof       = fact-ref | "exact" ;
instruction = "nop"
            | "drop", WS, fact-name, [ "..", fact-name ]
            | fact-name, "=", ("lin.le" | "lin.eq"), WS,
                expr-name, linear-list
            | fact-name, "=", "eq.pair", WS, expr-name, ",",
                fact-ref, ",", fact-ref
            | fact-name, "=", "box.le", WS, expr-name
            | fact-name, "=", "weaken.le", WS, expr-name, ",",
                fact-ref, ",", number
            | fact-name, "=", "int.aff", WS, expr-name, ",",
                number, linear-list
            | fact-name, "=", ("round.hi" | "round.lo"), WS,
                expr-name, ",", fact-ref, ",", fact-ref
            | "map.copy", WS, output-var, ",", input-var
            | "map.const", WS, output-var, ",", number
            | "map.aff", WS, output-var, ",", expr-name
            | ("check.row.hi" | "check.row.lo" | "check.row.eq"),
                WS, output-row, ",", proof
            | ("check.var.hi" | "check.var.lo" | "check.var.int"),
                WS, output-var, ",", proof
            | ("check.obj.eq" | "check.obj.le"), WS, proof
            | split ;
\end{lstlisting}

\begin{lstlisting}[caption={终止、split 与 case 语法}]
fact-list   = fact-ref, { ",", fact-ref } ;
output-list = fact-name, { ",", fact-name } ;
yield       = "yield", [ WS, fact-list ] ;
false       = "close.false", WS, fact-ref ;
branch-body = { instruction, NL }, (yield | false), NL ;
code-body   = { instruction, NL }, ("ret" | false), NL ;
split       = "(", [ output-list ], ")", "=", split-op ;
split-op    = "split.bin", WS, input-var, "{", NL,
                "case 0:", NL, branch-body,
                "case 1:", NL, branch-body, "}"
            | "split.int.thr", WS, fact-ref, ",", number, "{", NL,
                "case le:", NL, branch-body,
                "case ge:", NL, branch-body, "}"
            | "split.int.range", WS, fact-ref, ",", fact-ref,
                ",", fact-ref, ",", number, ",", number, "{", NL,
                range-cases, "}"
            | "split.sign", WS, expr-name, "{", NL,
                "case le:", NL, branch-body,
                "case ge:", NL, branch-body, "}" ;
range-cases = "case", WS, integer, ":", NL, branch-body,
              { "case", WS, integer, ":", NL, branch-body } ;
\end{lstlisting}

该 EBNF 规定结构；类型、范围、case 完整性、scope 可见性和终止可达性属于静态及执行语义检查。token 之间可插入水平空白，除显式 \code{WS} 外不强制数量。线性列表的右花括号后由外围语句换行结束。assembler 必须报告源码位置，不能通过自动补全缺失 case 或 obligation 使不完整程序成为有效证书。

\section{工具链接口草案}

\begin{longtable}{p{0.19\textwidth}p{0.73\textwidth}}
\toprule
工具 & 接口职责 \\
\midrule
\endhead
\code{rpa-as} & 解析文本，读取 canonical models 以检查索引和 exact 简写，展开 expressions，解析 FactId 和 scopes，输出 RPX。允许实现 liveness 优化，但不得改变命题或丢掉 obligation。\\
\code{rpa-dis} & 先验证 RPX 的结构，再输出 canonical RPA；不把成功反汇编解释为证明通过。\\
\code{rpa-objdump} & 显示 header、section directory、opcode counts、scope 深度及 fact liveness 等结构信息。结构统计与验证运行统计必须明确区分。\\
\code{rp-verify} & 读取两个 canonical models，绑定摘要和调用者要求，执行两个 proof tapes，成功后报告 VERIFIED 及 contract。\\
\code{mps-canon} & 将明确支持的 MPS 子集转换为 canonical model；不支持的输入特性必须显式拒绝。\\
\bottomrule
\end{longtable}

\begin{lstlisting}[caption={计划中的模型相关汇编与验证接口}]
mps-canon source.mps -o source.mcm
mps-canon target.mps -o target.mcm

rpa-as --source source.mcm --target target.mcm \
    proof.rpa -o proof.rpx

rp-verify --require opt-equivalent \
    source.mcm target.mcm proof.rpx

rpa-dis proof.rpx -o proof.canonical.rpa
rpa-objdump --sections --opcodes --liveness proof.rpx
\end{lstlisting}

以上命令展示接口契约而非可立即运行的实现。verifier 的成功输出至少应包含 VERIFIED、实际检查的 contract 和两个 model 摘要。instruction 数、derived fact 数、peak live facts 和耗时应说明其统计定义，特别区分 proof tape 解码、exact arithmetic 和模型加载成本。任何失败必须返回非零退出状态，并给出失败 domain、指令位置、错误类别和简短原因。

\section{扩展边界与保留语法}

基础 v0.1 不接受 \code{.batch} 或 \code{call.thm}。批量编码区 \code{0x80--0xBF} 和 \code{CALL\_THM} 的 \code{0xF0} 仅为未来扩展保留；基础 loader 遇到它们必须返回 \code{E_UNSUPPORTED}。

未来 batch 扩展必须定义完整的 core instruction 展开，包括 FactId 分配、scope、错误行为和 obligation 更新，不能通过批处理增加可证明命题。未来 theorem 扩展必须显式标识 trusted checker 和 TheoremId，不能用新名字绕过 trusted semantics 的版本审查。
